% TOP_PROBLEM: {"problemId": "problem.hadwiger-boltyanski-illumination-conjecture", "canonicalTitle": "Hadwiger-Boltyanski Illumination Conjecture", "releaseRank": 473, "primaryDomain": "combinatorics_discrete_geometry", "primaryDomainLabel": "Combinatorics & discrete geometry", "displayStatus": "open", "releaseStatus": "open"}
% !TeX program = lualatex
% Definition and sources only; this file is not an LLM solution attempt.
\documentclass[11pt]{article}
\usepackage[margin=25mm]{geometry}
\usepackage{fontspec}
\setmainfont{Latin Modern Roman}
\usepackage{amsmath,amssymb,mathtools}
\usepackage{unicode-math}
\setmathfont{Latin Modern Math}
\usepackage{xurl,hyperref}
\hypersetup{colorlinks=true,urlcolor=blue}
\setcounter{secnumdepth}{0}
\setlength{\parindent}{0pt}
\setlength{\parskip}{5pt}
\setlength{\emergencystretch}{3em}
\begin{document}
\section*{\texorpdfstring{Hadwiger-Boltyanski Illumination Conjecture}{Hadwiger-Boltyanski Illumination Conjecture}}\label{473}
\subsection{Definitions and mathematical statement}
A convex body $K\subset{\mathbb{R}}^n$ is compact and convex with nonempty interior. A direction $v\in S^{n-1}$ illuminates $x\in\partial K$ if $x+tv\in\operatorname{int}K$ for some $t>0$. Its illumination number is
\[
 I(K)=\min\{|D|:D\subset S^{n-1},\
 \forall x\in\partial K\ \exists v\in D\ \exists t>0:\
 x+tv\in\operatorname{int}K\}.
\]
The conjecture is $I(K)\le2^n$ for every dimension and every convex body. The equivalent external-point-source formulation lets a ray from a source continue through the boundary point into the interior. One direction may illuminate many boundary points; the finite collection must cover all of the boundary, not merely almost every point.
\par\smallskip\textit{Formulation and concept references:} \href{https://arxiv.org/pdf/2507.08712}{[S1]}.

\subsection{Short English statement}
Can the boundary of every n-dimensional convex body be illuminated using at most two to the n light directions?
\subsection{Sources}
\begin{itemize}
\item[S1]Arman, A., Kaire, J.S., Prymak, A., 'Illumination number of 3-dimensional cap bodies', Discrete Math. 349(7) (2026); arXiv:2507.08712.
\url{https://arxiv.org/pdf/2507.08712}.
\textit{Formulation source.}
\end{itemize}
\end{document}
